% 全书导读（前言部分，不编号）
\chapter{导读}

本笔记按原讲义的五章加一个附录组织，依据课程讲义 \emph{Convex analysis and optimisation}（2023 年 9 月 28 日版）整理成中文。
原讲义的错漏已在正文中直接改正，逐条记录见同目录下的《修改报告.md》。

\paragraph{工作空间。}
讲义的大部分结论在赋范空间中成立，读者可以始终把“（局部凸）拓扑向量空间”换成“赋范空间”甚至 $\R^d$ 来理解。
只有在讨论弱拓扑、弱*拓扑时才真正需要一般的拓扑向量空间，相关预备知识放在附录~\ref{ch:appA}。

\paragraph{各章内容。}
\hyperref[ch:ch01]{第 1 章}讨论凸集：Hahn--Banach 定理、凸集的分离定理、闭凸集是闭半空间之交，以及支撑函数与法锥。
\hyperref[ch:ch02]{第 2 章}讨论凸函数：连续性与局部 Lipschitz 性、下半连续凸函数、凸共轭与 Fenchel--Moreau 定理，以及极小点的存在性。
\hyperref[ch:ch03]{第 3 章}讨论次微分：方向导数、Gâteaux 与 Fréchet 可微、次微分的运算法则、最优性条件，以及凸函数几乎处处可微。
\hyperref[ch:ch04]{第 4 章}讨论邻近算子与邻近点算法，以及共轭函数的邻近算子（Moreau 分解）。
\hyperref[ch:ch05]{第 5 章}讨论凸对偶：扰动框架，以及它的三个应用，即 Lagrange 对偶、Fenchel--Rockafellar 定理和最优输运中的对偶。

\paragraph{主线。}
全课程的核心工具是“闭凸集等于包含它的闭半空间之交”（第 1 章），它来自 Hahn--Banach 定理。
把这个结论用到上图上，就得到“下半连续凸函数等于其仿射下界的上确界”，即 Fenchel--Moreau 定理 $f^{**}=f$（第 2 章）；
用到支撑超平面上，就得到次梯度的存在性与次微分的运算法则（第 3 章）；
用到扰动函数的值函数上，就得到凸对偶理论（第 5 章）。
第 4 章的邻近算子则是把这些理论变成算法的桥梁。
